\section{Yet to be implemented}\label{sec:incomplete}

Every expository chapter of this document has marked its unbuilt edges in place --- a clause
here, a callout there --- and deferred the consolidated view to this chapter. This is that
view: the complete inventory, at pinned commit \code{254cd41}, of what QtRocket has built
but not wired in, what it has wired in with named gaps, and what it has not started. The
companion Architecture Review (\srcf{docs/ArchitectureReviewLatex/}) calls the first
category ``a ring of machinery that is built, unit-tested, and consumed by nothing,'' and
that ring is the defining feature of the project's frontier: QtRocket builds and tests its
seams before their consumers exist, so most of what is missing is not missing code --- it is
a missing caller. The tree says so itself, in the aerodynamics suite's own words:
``nothing in production consumes this in P2 -- these tests pin the seam for P5''
\src{model/tests/AeroTests.cpp}{18}.

This chapter is written as a map for the next engineer, not as a critique. Each entry
records three things: what exists, where it lives, and what consuming it would require.
The ordering within each category is roughly by size of the waiting machinery; the
\emph{ranking} by user impact belongs to the Architecture Review and is reproduced, with
attribution, in \cref{sec:incomplete-sequencing}.

\begin{designrule}[The status vocabulary, and how absence is grounded]
Three status words partition this inventory, with the companion review's exact senses.
\emph{Dormant}: built and tested, with zero production consumers --- unreachable from any
shipping code path. \emph{Partial}: reachable from a production path, but with named gaps.
\emph{Absent}: no trace in the tree --- and every absence claim below is grounded in a
failed search of the source, not an impression. The distinction is a planning tool, not
taxonomy for its own sake: consuming a dormant seam is weeks of work where building an
absent subsystem is months.
\end{designrule}

\subsection{Dormant seams}\label{sec:incomplete-dormant}

Seven seams sit finished, most of them unit-tested, with no production caller. Each names, implicitly or in
an in-source comment, the work it is waiting for.

\begin{description}

\item[The Barrowman layer and \code{aeroData}.] The largest dormant seam is the entire
center-of-pressure (CP) pipeline. Per-part Barrowman contributions exist for the cone and
the fin set, expressed as a normal-force slope \code{cnAlpha} and a slope-weighted CP moment
\code{cnAlphaXcp} \src{sim/Aero.h}{32}; \code{Part::getCompositeAero()} folds them
additively over resolved placement onto the nose-tip datum \src{model/parts/Part.cpp}{231},
the same datum as the composite center of gravity (CG), so the static margin is literally
\cppinline|cp() - cg()| --- a subtraction by construction, and the seam's own documentation
states that intent \src{sim/Aero.h}{25}. All of it is unit-tested
(\srcf{model/tests/AeroTests.cpp}, \srcf{model/tests/NoseConeTests.cpp}) and none of it is
called from production: the only callers of \code{getCompositeAero} are three test files
\src{tests/PlacementInvarianceTests.cpp}{106}, and the \code{aeroData} member staged for the
result on the model--sim bridge is default-constructed and never read
\src{model/Propagatable.h}{61}, kept compiling only by a back-compatibility alias
\src{sim/Aero.h}{66}. Consuming the seam requires an angle-of-attack computation (none
exists anywhere in the tree), a normal-force or restoring-moment term in the force path,
and a display surface; the cheapest first consumer is a live CP/CG/static-margin readout,
which needs no new physics at all. Two latent defects gate that consumption and are covered
with the 6-DOF cluster below (\cref{sec:incomplete-sixdof}); the full analysis of the layer
is \cref{sec:aero}'s.

\item[The wind model.] \code{WindModel} is the purest specimen of the pattern: compiled
into the \code{sim} library \src{sim/CMakeLists.txt}{24}, its entire behavior is
\cppinline|return Vector3(0.0, 0.0, 0.0);| \src{sim/WindModel.cpp}{18}, and a repo-wide
search finds no instantiation, no call site, and no slot for it in \code{Environment}
(\srcf{sim/Environment.h} holds gravity, atmosphere, and geoid members only). The drag law
consumes ground-relative velocity directly \src{model/RocketModel.cpp}{88}, so airspeed
equals ground speed by construction. Consuming it requires an \code{Environment} slot and a
relative-velocity substitution in \code{getForces()} --- but, as the review observes, wind
is only worth plumbing after attitude exists: a point mass cannot weathercock, so wind
would shift drag without producing any behavior wind exists to model.

\item[Ignition offsets and ejection delays.] The motor subsystem carries two time-related
seams. First, arbitrary ignition times: \code{MotorModel::startMotor(double startTime)}
accepts any start time \src{model/MotorModel.h}{328}, and the mass model honors the offset
by evaluating at $\mathtt{simTime} - \mathtt{ignitionTime}$ \src{model/MotorModel.cpp}{36}
--- but \code{getThrust} passes absolute simulation time through to the curve
\src{model/MotorModel.cpp}{82}, and \code{ThrustCurve::setIgnitionTime}, the piece that
would shift the thrust samples, has zero call sites, with its \code{maxTime} adjustment
commented out inside the definition \src{model/ThrustCurve.cpp}{44}. The asymmetry is
masked because the only production ignition is \cppinline|startMotor(0.0)| at launch
\src{model/RocketModel.cpp}{108}; any future air-start would burn propellant on one clock
and produce thrust on another until the two paths are unified (\cref{sec:motors}). Second,
ejection delays: both file loaders parse them --- the RockSim loader splits the CSV list
\src{model/RSEDatabaseLoader.cpp}{59}, the RASP loader splits on \code{-} with 1000 as the
plugged sentinel \src{model/RASPLoader.cpp}{119}, and the thrustcurve.org client leaves them
an explicit TODO \src{model/ThrustCurveClient.cpp}{198} --- they are stored
\src{model/MotorModel.h}{303} and round-tripped through the motor database XML, and no
simulation, GUI, or CLI code reads them. Their consumer is the absent recovery-event system
(\cref{sec:incomplete-absent}).

\item[\code{deriveReferenceAreaFromGeometry()}.] The geometry-derived reference area is
built end to end --- \code{RocketModel} exposes it \src{model/RocketModel.cpp}{48} over
\code{Part::maxFrontalReferenceArea()}'s widest-frontal-disc convention
\src{model/parts/Part.cpp}{251} --- and has no production caller. The header's promise that
the geometry default applies whenever the user has not overridden the area
\src{model/RocketModel.h}{76} is implemented by no code path: design load applies the saved
area only when \code{referenceAreaOverridden="true"} \src{model/DesignSerializer.cpp}{265},
and all 25 bundled fixtures say \code{"false"}, so every loaded design flies on whatever
scalar was last staged --- by default the 38\,mm disc. Drag therefore does not scale with
airframe class; \cref{sec:aero} carries that verified defect in full. Consumption here is
one call: wire the derivation into the non-overridden branch, and the seam's convention
(fins deliberately report the body disc, not tip extent, \src{model/parts/FinSet.h}{71})
already matches standard practice.

\item[\code{radialSeamCheck} (the Layer-1 diagnostic).] The placement design specifies three
diagnostic layers; the shipped gate wires in only the second, the envelope sweep
(\cref{sec:placement}). The first --- a per-seam, $O(1)$ radius-compatibility check
dispatched on seat kind --- is implemented \src{model/parts/Placement.cpp}{56} and
unit-tested \src{model/tests/ResolverSweepTests.cpp}{194}, and nothing in production calls
it: the placement cache's re-resolve runs the sweep alone \src{model/parts/Part.cpp}{170}.
The function even documents its own missing caller --- its \code{zWorld} field is left
parent-local behind the comment ``the root-frame z is added once a pose is known''
\src{model/parts/Placement.cpp}{87}, a promise no caller exists to fulfill. Consuming it
means invoking it per attachment during resolve, converting a wrong-radius mate from a
sweep-time overlap (or a silent pass) into an immediate, located seam diagnostic.

\item[Per-sample RockSim mass and CG histories.] The bundled motor file ships
manufacturer-measured mass and CG values on every thrust sample
\src{data/Aerotech.rse}{19}, and the loader reads only \code{t} and \code{f}
\src{model/RSEDatabaseLoader.cpp}{91}. QtRocket instead always reconstructs burn mass with
the 128-point constant-$I_{sp}$ depletion curve \src{model/MotorModel.cpp}{132} and holds
the motor's own center of mass fixed \src{model/parts/Motor.cpp}{12} --- the fallback a
mature tool uses when a file carries no measured data, applied unconditionally. Consuming
the discarded columns requires a per-motor measured mass curve in place of the
reconstruction and intra-motor CG travel during the burn; the latter also depends on motor
placement becoming real (\cref{sec:incomplete-partial}), since a CG history acting at the
airframe CM would move the right mass at the wrong station.

\item[The \code{isp} member.] The smallest seam, one line: specific impulse is computed
whenever motor metadata is installed, $I_{sp} = I_{tot} / (g_0 \cdot m_{prop})$
\src{model/MotorModel.cpp}{130}, stored \src{model/MotorModel.h}{340}, and read by nothing.
It is the natural input for a future propellant-performance display or a
mass-flow-consistency check on ingested files; today it is dead weight computed on every
metadata write.

\end{description}

\subsection{The 6-DOF cluster}\label{sec:incomplete-sixdof}

The largest coordinated group of dormant machinery is the scaffolding for six degrees of
freedom (6-DOF) --- rotational state, torques, and the inertia pipeline that would drive
them. Five pieces exist, typed and in place, and none is consumed:

\begin{itemize}
\item \textbf{The torque interface.} \code{getTorques(t)} is a pure virtual on the
  model--sim bridge \src{model/Propagatable.h}{30}; \code{RocketModel}'s override returns
  the zero vector unconditionally \src{model/RocketModel.cpp}{94}. A repo-wide search finds
  no call site in \srcf{sim/}, \srcf{gui/}, \srcf{cli/}, or the controller --- only the
  declaration, the zero override, and a test mock \src{tests/PropagatorTests.cpp}{43}.
\item \textbf{The inertia accessor.} \code{getCompositeInertiaTensor(t)} is declared
  \src{model/Propagatable.h}{33} and implemented as a pass-through to the part tree
  \src{model/RocketModel.cpp}{43}, and no production code calls the virtual; the per-step
  recording reaches the tree directly through \code{writeMassProperties}
  \src{model/RocketModel.cpp}{59}. The composite $I(t)$ and CG$(t)$ that
  \cref{sec:parts} shows the tree computing so carefully have exactly one time-varying
  consumer today: the per-step state snapshot (the CLI's \code{listparts} also prints a
  static $t=0$ composite).
\item \textbf{The quaternion state slots.} \code{StateData} carries \code{orientation} and
  \code{orientationRate} quaternions, a direction-cosine matrix, and 3-2-1 Euler angles
  \src{sim/StateData.h}{46}; no code anywhere writes or reads them. Both quaternions are
  initialized to the all-zero quaternion --- not the identity rotation, not any rotation
  --- and the members' ``(vector, scalar)'' annotation sits one constructor-order hazard
  away from Eigen's $(w, x, y, z)$ value constructor (\cref{sec:simcore}). The first
  consumer inherits an invalid orientation plus a documented convention trap.
\item \textbf{The orientation integrator.} The rotational half of the propagator exists as
  a commented-out member whose interface is already worked out
  (\cref{lst:incomplete-orientation}): the same \code{DESolver} machinery, instantiated for
  quaternions, with a time-keyed callback so torques are sampled at each stage's node time
  exactly as forces are. The solver templates admit the quaternion instantiation by
  \cppinline|static_assert| \src{sim/RK45Solver.h}{51}, though only \code{Vector3}
  instantiations exist \src{sim/Integrator.h}{57}, and the stage arithmetic would not
  compile for \code{Eigen::Quaterniond} without adaptation --- the template is a staged
  intention, not working rotational machinery.
\item \textbf{The recorded seam.} Every accepted step snapshots mass, CG, and the full
  inertia tensor into the state history \src{sim/Propagator.cpp}{132}, and the header names
  the purpose: these are ``the seam the 6-DOF rotational ODE will read''
  \src{sim/StateData.h}{59}. On the placement side, every \code{StationLink} carries a
  \code{childRot} quaternion reserved for 6-DOF, identity in 3-DOF
  \src{model/parts/Placement.h}{55}, deliberately left out of the design file format while
  it is identity by construction \src{model/DesignSerializer.cpp}{85}.
\end{itemize}

\begin{listing}[htbp]
\begin{minted}{cpp}
   std::unique_ptr<sim::Integrator> linearIntegrator;
   // The 6-DOF orientation integrator will use the same DESolver<Quaternion> interface, with a
   // time-keyed ODE callback so getTorques() can be sampled at each stage's node time like getForces().
//   std::unique_ptr<sim::RK4Solver<Quaternion>> orientationIntegrator;
\end{minted}
\caption{The rotational half of the propagator, present as a commented-out member with its
design already stated in place \srccap{sim/Propagator.h}{100}. The dormant seam names its
own consumer: a quaternion ODE fed by \texttt{getTorques()} at stage node times.}
\label{lst:incomplete-orientation}
\end{listing}

\begin{keyinsight}[Two defects gate 6-DOF consumption]
The Architecture Review names two latent defects that must be fixed \emph{before} this
cluster gains a consumer, because both leave the 3-DOF trajectory untouched and become
flight errors the moment rotation is integrated. First, the fin set's chordwise frame is mirrored: its mass
and CP math use aft-positive from-root-leading-edge distances directly as $+z$ offsets
\src{model/parts/FinSet.cpp}{57}, placing the composite fin CP aft of truth and the fin-set
mass forward of truth, and the unit tests pin the mirrored values as regression anchors
\src{model/tests/FinSetTests.cpp}{207} --- the suite enforces the defect rather than
detecting it (\cref{sec:aero}). Second, the composite-inertia walk deliberately omits the
$R\,I\,R^{\mathsf{T}}$ child-tensor rotation --- identity in 3-DOF, skipped ``to stay
bit-stable'' --- and the in-source TODO already prescribes the regression test that must
accompany its addition \src{model/parts/Part.cpp}{217}. Fixed in that order, the fin-set
mirror is flushed out while it is still a display bug rather than a flight bug.
\end{keyinsight}

\subsection{Partial areas}\label{sec:incomplete-partial}

Six areas are reachable from production paths today, each with a named gap.

\begin{description}

\item[Placement authoring.] The model and the serializer support the full placement
vocabulary --- three seat kinds \src{model/parts/Placement.h}{32}, arbitrary fractional
stations, and gaps, all round-tripped through the design file
\src{model/DesignSerializer.cpp}{54} --- but the only authoring surface, the REPL's
\code{addpart}, attaches everything with \cppinline|abut(offset.z())|
\src{cli/Repl.cpp}{872}, and parses \code{x=}/\code{y=} offsets only to discard them
\src{cli/Repl.cpp}{189}. \code{Abut} is the only seat kind reachable from any front end;
\code{NestInBore} and \code{OnSurface} placements exist today only in hand-edited
\code{.qrd} files and C++ test fixtures. The file format is ahead of the authoring
vocabulary (\cref{sec:placement}).

\item[The GUI design surface.] At the pinned commit the main window's design widgets are
constructor-only stubs --- \code{RocketTreeView} is an empty constructor
\src{gui/RocketTreeView.cpp}{3}, \code{RocketModelerView} the same
\src{gui/RocketModelerView.cpp}{3} --- and the File menu's New/Open/Save/Save~As/Close
actions ship \code{enabled="false"} \src{gui/MainWindow.ui}{117}; exactly three actions are
wired: Quit, Save Motor Database, and About \src{gui/MainWindow.cpp}{42}. Every capability
the widgets promise exists one layer down in the CLI (\cref{sec:interfaces}). Two commits
on the unmerged \code{RocketTreeViewImplementation} branch add a read-only part tree and
\code{.qrd} open/save wiring; that work is in progress and is not reviewed here.

\item[GUI responsiveness and change notification.] Everything the GUI does, it does on the
Qt event-loop thread: the launch slot flies the whole simulation synchronously inside the
button handler \src{gui/CannonballTab.cpp}{116}, and one motor-search click can issue up to
$1 + 100$ blocking HTTP round trips (\cref{sec:interfaces}) --- there is no \code{QThread},
\code{QtConcurrent}, or \code{QNetworkAccessManager} anywhere under \srcf{gui/}. In the
other direction the model cannot tell the GUI anything: no structure-change notification
out of \code{RocketModel} exists at the pinned commit (a failed repo-wide search for any
callback or signal), so a future part-tree widget has nothing to subscribe to; the unmerged
\code{RocketTreeViewImplementation} branch prototypes such a callback.

\item[Analysis and plotting.] The GUI's analysis window draws exactly three fixed graphs
--- altitude \src{gui/AnalysisWindow.cpp}{49}, vertical velocity
\src{gui/AnalysisWindow.cpp}{74}, and the motor thrust curve
\src{gui/AnalysisWindow.cpp}{86} --- and never surfaces the flight statistics or the typed
termination reason the controller exposes \src{QtRocket.h}{55}: a user whose run aborted
for no-liftoff gets an altitude plot and no explanation. The CLI prints the statistics and
writes a fixed fourteen-column CSV \src{cli/Repl.cpp}{104}; there is no configurable
plotting, no event annotation, and no stability output on any surface.

\item[The motor database ships empty.] Exactly one motor file is bundled,
\srcf{data/Aerotech.rse}, and only the test binaries can find it, through the
\code{QTROCKET\_DATA\_DIR} compile definition \src{tests/CMakeLists.txt}{16}; neither front
end loads any motor data at startup (a failed search --- no production reference to the
data directory exists). Every GUI or CLI session begins with an empty database and a manual
import (\cref{sec:motors}).

\item[Single motor, single stage.] One motor is a structural assumption: the model borrows
a handle to the \emph{first} \code{Motor} node found by depth-first search
\src{model/RocketModel.cpp}{13}, a second motor is silently ignored for thrust while still
contributing mass, and \code{Stage.h} is a commented-out include under the comment ``Not
yet'' \src{model/RocketModel.h}{22}. The one motor is also pinned to the wrong place: it is
attached CM-to-CM with the airframe under an offset documented ``Zero today'' with no
setter \src{model/RocketModel.h}{131}, so the CG-travel-during-burn machinery of
\cref{sec:parts} acts at the airframe CM rather than in an aft mount, and the motor's
placement link is never persisted \src{model/DesignSerializer.cpp}{114} --- lossless today
only because the offset cannot move.

\end{description}

\subsection{Absent subsystems}\label{sec:incomplete-absent}

The following have no trace in the tree. Each claim rests on a failed search of the source
at the pinned commit, not on an impression.

\begin{description}

\item[Recovery and flight events.] Nothing changes at apogee. Once off the pad, the only
phase logic in the simulator is the terminate predicate --- descending below the launch site
\src{model/RocketModel.cpp}{62} --- so every flight ends ballistically, descending under
the ascent's drag law; apogee exists only as a running maximum in the statistics
\src{sim/TrajectoryStatistics.h}{29}. There is no deployment event, no parachute or
streamer model, no drag switch at any altitude. The ejection delays every loader ingests
(\cref{sec:incomplete-dormant}) are this subsystem's staged input, and the review points at
the typed termination taxonomy \src{sim/Propagator.h}{35} as the natural seed of an event
vocabulary --- liftoff, burnout, apogee, ejection, touchdown.

\item[Drag build-up, Mach, and Reynolds.] The per-part drag field exists and every producer
returns zero: \code{cd} stays 0 on the body tube --- ``skin friction over getWettedArea()
is its eventual contribution'' \src{model/parts/BodyTube.cpp}{41} --- on the nose cone
\src{model/parts/ConicalNoseCone.cpp}{67}, and on the fin set
\src{model/parts/FinSet.cpp}{89}; the wetted-area getters staged for the purpose have no
production consumer \src{model/parts/BodyTube.h}{51}. Above them, the atmosphere interface exposes
speed of sound and dynamic viscosity \src{sim/AtmosphericModel.h}{17}, and the US Standard
Atmosphere implements both correctly \src{sim/USStandardAtmosphere.cpp}{143} --- those two
overrides have no caller of any kind, production or test. No Mach number or Reynolds number
is computed anywhere (failed search); the single user-set $C_d$ is speed-independent, so an
upper-class flight passes Mach~1 under subsonic-constant drag without comment
(\cref{sec:aero,sec:environment}).

\item[Staging, clusters, and flight configurations.] No \code{Stage} type exists anywhere
under \srcf{model/} (the commented include above is the only mention); there is no motor
cluster support --- the single-motor DFS makes a second motor physically inconsistent
rather than illegal --- and no concept of named flight configurations pairing a design with
alternative motors.

\item[Launch rod, launch site, and Coriolis.] The only pad modeling is a hold clamp ---
the rocket is pinned at rest at $z = 0$ until the first powered step lifts it
\src{sim/Propagator.cpp}{111} --- and a no-liftoff abort \src{sim/Propagator.cpp}{146}.
There is no guided rod phase, rod length, or departure-velocity constraint; since attitude
does not exist, a rod could not constrain the thrust direction anyway. There is likewise no
launch-site concept: the geoid ignores its latitude/longitude arguments
\src{sim/SphericalGeoidModel.cpp}{18}, the spherical gravity model queries it once with a
literal $(0, 0)$ placeholder \src{sim/SphericalGravityModel.cpp}{20}, altitude above the
pad is implicitly altitude above sea level in every atmosphere lookup, and no Coriolis or
Earth-rotation term appears anywhere (failed search).

\item[Foreign design formats.] QtRocket imports foreign \emph{motor} data through three
paths (\cref{sec:motors}) and foreign \emph{design} data through none: there is no
OpenRocket \code{.ork} reader or writer, no RockSim \code{.rkt} design import, and no
export beyond the trajectory CSV \src{cli/Repl.cpp}{104}. Its own legacy 0.1 design files
are rejected fail-closed with ``must be re-created'' rather than migrated
\src{model/DesignSerializer.cpp}{194}. Combined with the five-type part catalog
(\cref{sec:parts}), the \code{.qrd} format is a closed world at this commit
(\cref{sec:persistence,app:or-features}).

\item[Optimization and dispersion.] No optimizer, no Monte-Carlo or dispersion analysis;
searches for either return nothing. The deterministic core and the scriptable REPL make
external parameter sweeps feasible today --- the design-matrix suite is exactly such a
sweep run as a test (\cref{sec:testing}) --- but nothing in the product performs one.

\item[Packaging and internationalization.] Distribution is a single install rule covering
the GUI target only \src{CMakeLists.txt}{308}; no CPack configuration exists (failed
search), the macOS bundle identifier is the Qt template placeholder
\src{CMakeLists.txt}{301}, and the CLI, the visualizer, and the motor data are not
installed at all. Internationalization is dead scaffolding: the build generates translation
catalogs nothing references \src{CMakeLists.txt}{201}, and the GUI loads a translator from
a resource prefix that is never populated \src{gui/GuiRunner.cpp}{40}, under a TODO
admitting English-only support \src{gui/GuiRunner.cpp}{34}.

\end{description}

\subsection{A suggested order of work}\label{sec:incomplete-sequencing}

Ranking this inventory is deliberately not this document's job; the Architecture Review
does it, and its recommendation is reproduced here because it follows directly from the
dormant-versus-absent distinction above. The review's sequencing
(\srcf{docs/ArchitectureReviewLatex/sections/11-roadmap.tex}) wires seams to consumers in
dependency order:

\begin{enumerate}
\item \textbf{Consume the Barrowman seam.} Fix the fin-set chordwise mirror, add the
  $R\,I\,R^{\mathsf{T}}$ rotation with the pinned-identity regression test the TODO already
  prescribes \src{model/parts/Part.cpp}{220}, and give
  \code{deriveReferenceAreaFromGeometry()} its caller. The payoff is a CP/CG/static-margin
  readout in the CLI and GUI --- no new physics, the first production exercise of
  \code{getCompositeAero}, and the sign defect flushed while it is still a display bug
  rather than a flight bug.
\item \textbf{Recovery and flight events.} Independent of rotation: grow the termination
  taxonomy \src{sim/Propagator.h}{35} into an event vocabulary, consume the ejection
  delays, and switch drag at deployment. This replaces the ballistic terminate with a full
  flight --- and a minutes-long parachute descent forces the simulation off the GUI thread,
  which in turn makes the controller's initialization race (\cref{sec:overview}) worth
  fixing first.
\item \textbf{Six degrees of freedom, then wind.} Reinstate the orientation integrator
  (\cref{lst:incomplete-orientation}), fixing the quaternion zero-initialization and the
  constructor-order hazard before any consumer reads them; align thrust with the body axis;
  derive restoring torques from step~1's aero profile. Only then plumb the wind model ---
  a point mass cannot weathercock, so connecting it earlier would shift drag without
  producing any behavior wind exists to model.
\item \textbf{The GUI design editor.} Pure consumption of a proven core: the CLI already
  demonstrates the full author-persist-fly loop against the same API, the 25 committed
  \code{.qrd} fixtures are a ready test corpus, and the in-flight branch supplies the tree
  view. What remains is the property panel, part add/remove/edit, and closing the CLI's
  seat-and-station authoring gap.
\end{enumerate}

\begin{designrule}[Consumption before construction (the review's verdict)]
``Almost nothing in Tiers 1--2 requires new architecture. The aero profile, the torque
interface, the quaternion state, the ejection delays, the termination taxonomy, the
authoring API, and the serializer all exist and are tested. The highest-leverage sequencing
wires existing seams to consumers in dependency order, fixing the two known latent defects
before their numbers become user-visible''
(\srcf{docs/ArchitectureReviewLatex/sections/11-roadmap.tex}).
\end{designrule}

Staging, clusters, the multi-host joint graph, and packaging stay off this critical path by
design: none blocks the four steps, and each gets cheaper after them --- staging needs the
event system, clusters need the joint graph.

\subsection{The placement system's forward hook}\label{sec:incomplete-hook}

The Part Placement whitepaper (\srcf{docs/PartPlacementDesignLatex/}) closes with the same
forward posture this chapter documents, and one of its designed pieces is still unbuilt:
Layer~3 of its diagnostics ladder, the degrees-of-freedom accounting. In 3-DOF a single
axial \code{StationLink} fully pins a child; under 6-DOF one link leaves a rotational
degree of freedom unpinned, and the design specifies a cheap detector that warns when a
child has only one link under free orientation
(\srcf{docs/PartPlacementDesignLatex/sections/06-diagnostics.tex}). No such code exists in
\srcf{model/parts/} --- verified by reading \srcf{model/parts/Placement.cpp} and
\srcf{model/parts/Part.cpp} end to end --- so the detector is designed, not shipped:
absent, with its specification already written.

The rest of the placement system's 6-DOF path is the inventory's most finished forward
hook. The schema is already shaped for rotation --- \code{childRot} rides in every stored
link as an identity quaternion \src{model/parts/Placement.h}{55}, \code{Pose::compose()} is
the full rotate-then-translate rigid-transform law that degenerates to vector addition
under identity \src{model/parts/Placement.h}{81} --- so the design document reduces 6-DOF
placement to three localized edits: stop forcing \code{childRot} to identity at authoring
and load time, add the one $R\,I\,R^{\mathsf{T}}$ line to the composite-inertia pass, and
feed the integrated attitude into the resolver's root pose
(\srcf{docs/PartPlacementDesignLatex/sections/09-sixdof.tex}). Zero schema change, zero
signature change, no second resolver.

That is the shape of this whole chapter, and the review closes with the image that fits it:
the hard structural machinery --- composite inertia, placement resolution, Barrowman
composition, adaptive integration, the termination taxonomy --- is built, tested, and
waiting behind seams that name their consumers. In the review's words, ``the rooms are
framed and inspected; the remaining work is opening the doors between them''
(\srcf{docs/ArchitectureReviewLatex/sections/11-roadmap.tex}).
